\begin{frame}{Frequência de amostragem}
    A frequência de amostragem indica quantas medições o ADC realiza por segundo.

    \[
        f_s = \text{número de amostras por segundo}
    \]

    Para representar adequadamente um sinal, utiliza-se como referência o
    teorema de Nyquist:

    \[
        f_s \geq 2 \times f_{\text{máx}}
    \]

    \begin{exampleblock}{Exemplo}
        Para medir um sinal de até $1$ kHz, a frequência de amostragem deve ser,
        no mínimo, aproximadamente $2$ kHz.
    \end{exampleblock}
\end{frame}